\SourceLeaf{319}%
bis zur fünften Ordnung ausschließlich richtig sind, wenn man \(s\) als Größe erster Ordnung betrachtet. Ebenso sieht man, dass man in den eingeklammerten Gliedern auf der rechten Seite \(q,\alpha l\) und \(u\) ohne Genauigkeitsverlust durch \(s\mathord{.}\cos U\), \(s\mathord{.}\frac{\sin U}{\cos Q}\), \(s\mathord{.}\sin U\mathord{.}\operatorname{tang}Q\) ersetzen darf.

\Article{22}
Zunächst müssen die Größen \(B,b,T,t,r\), die sich auf die Ellipsoidoberfläche beziehen, mit ihren Entsprechungen auf der Kugeloberfläche \(Q,q,U,u,As\) verglichen werden. Alle Gleichungen, die dafür hier aufgestellt werden, sind mindestens bis zur dritten Ordnung einschließlich genau. Dass diese Bedingung erfüllt ist, lässt sich aus der Herleitung selbst leicht erkennen.

Wendet man die in Artikel 8 angegebene Reihe auf unsere beiden Punkte an, so sind die dort allgemein mit \(p\) und \(q\) bezeichneten Größen in unserer jetzigen Schreibweise auszudrücken
\begin{center}
\begin{tabular}{ll}
für den ersten Punkt durch & \(B+\frac12b-P\) und \(\frac12q\)\\
für den zweiten Punkt durch & \(B-\frac12b-P\) und \(-\frac12q\)
\end{tabular}
\end{center}
und damit haben wir die beiden Gleichungen
\begingroup\fontsize{10}{12}\selectfont
\[
\begin{aligned}
B+\tfrac12b={}&P+\frac{\cos\theta}{2\cos\varphi}\mathord{.}q-\frac{3ee}{8\cos\varphi^2}\mathord{.}\cos P\mathord{.}\sin P\mathord{.}qq\\
&+\frac{ee}{16\cos\varphi^3\cos\theta}(-\cos P^2+\sin P^2+ee(5\cos P^2\mathord{.}\sin P^2-\sin P^4))q^3\\
B-\tfrac12b={}&P-\frac{\cos\theta}{2\cos\varphi}\mathord{.}q-\frac{3ee}{8\cos\varphi^2}\mathord{.}\cos P\mathord{.}\sin P\mathord{.}qq\\
&-\frac{ee}{16\cos\varphi^3\cos\theta}(-\cos P^2+\sin P^2+ee(5\cos P^2\mathord{.}\sin P^2-\sin P^4))q^3
\end{aligned}
\]
\endgroup
Durch Addition und Subtraktion ergibt sich also
\[
\begin{aligned}
B={}&P-\frac{3ee}{8\cos\varphi^2}\mathord{.}\cos P\mathord{.}\sin P\mathord{.}qq\quad .\ .\ .\ .\ .\ .\ .\ .\ .\ .\ .\ .\ .\ .\ .\ .\ .\ .\ .\quad(4)\\
b={}&\frac{\cos\theta}{\cos\varphi}\mathord{.}q+\frac{ee}{8\cos\varphi^3\cos\theta}(-\cos P^2+\sin P^2\\
&\qquad+ee(5\cos P^2\mathord{.}\sin P^2-\sin P^4))q^3\qquad(5)
\end{aligned}
\]
Man sieht übrigens leicht, dass Gleichung (4) um Größen vierter Ordnung ungenau ist, Gleichung (5) dagegen nur um Größen fünfter Ordnung.

Um \(T\) und \(t\) mit \(U\) und \(u\) zu vergleichen, muss man die am Schluss von Artikel 15 hergeleiteten Formeln verwenden. Ihnen lag eine Voraussetzung
